\section{Supplementary Material}
\label{sec:supplement}

This supplement provides the full proofs of Theorems~\ref{thm:hessian}
and~\ref{thm:twoscale}, additional experiments referenced in the main
text, and robustness checks.

\subsection{Full proof of Theorem~\ref{thm:hessian}}
\label{sec:proof_thm1}

We restate Theorem~\ref{thm:hessian} for reference:

\begin{quote}
\emph{Under Assumptions \ref{ass:linear}--\ref{ass:inputlip}, the
effective condition number of the joint Gauss-Newton matrix satisfies
$\kappa_{\mathrm{eff}}(G) \geq c_1 \cdot \sqrt{\Cg/\Cf} - c_2$ for
some $c_1, c_2 > 0$ depending on the data covariance, the bottleneck
dimension $K$, and the model class of $g_\phip$.}
\end{quote}

The proof is structured as three lemmas: (i) the spatial block's top
eigenvalue scales with $\sqrt{\Cg}$ (Lemma~\ref{lem:phi_scaling},
Karakida-style mean-field FIM analysis); (ii) Cauchy interlacing
transfers this scaling to the full Hessian
(Lemma~\ref{lem:interlacing}); (iii) a residual-dimension rank bound
caps the smallest positive eigenvalue by a $\Cg$-independent constant
(Lemma~\ref{lem:schur}). We prove each in turn, then combine.

\subsubsection{Proof of Lemma~\ref{lem:phi_scaling}}

We show $\lambda_{\max}(J_\phip^\top J_\phip) = \Omega(\Cg)$ in
expectation over random Gaussian initialization, under
Assumptions~\ref{ass:reg} (twice-differentiable, $L$-Lipschitz
gradients) and \ref{ass:subg} (sub-Gaussian inputs with bounded
covariance condition number).

\paragraph{Step 1: Bounded per-pixel residual at initialization.}
Under Kaiming initialization, each pre-activation in the spatial
module $g_\phip$ has $\mathcal{O}(1)$ variance. The output logits
$\hat{y} = g_\phip(f_\thetap(X))$ are therefore distributed as
$\hat{y} \sim \mathcal{N}(0, \sigma_y^2 I_{N_{cls}})$ at initialization
in the mean-field limit, with $\sigma_y^2 = \Theta(1)$ independent
of $\Cg$. The softmax probabilities $p = \mathrm{softmax}(\hat{y})$
are concentrated near uniform $1/N_{cls}$, and the per-pixel residual
$r = p - y_{onehot}$ satisfies
\[
    \E\!\left[ \norm{r}^2 \right]
    \;=\; \frac{N_{cls}-1}{N_{cls}} \;+\; \mathcal{O}(\sigma_y^2)
    \;=\; \Theta(1).
\]
This bound is independent of $\Cg$.

\paragraph{Step 2: Trace of the Fisher Information block.}
Karakida et al.~\cite{karakida2019universal}, Theorem~3.1, give that
for a deep network at random Kaiming initialization, the expected
trace of the per-layer FIM block scales linearly with the layer's
parameter count:
\[
    \E\!\left[ \mathrm{tr}(F_\phip) \right]
    \;=\; \sum_{\ell=1}^{L} q_\ell \cdot \nu_\ell \cdot M_\ell,
\]
where $M_\ell$ is the parameter count of layer $\ell$, $q_\ell$ is
the squared activation magnitude at layer $\ell$, and $\nu_\ell$ is
the squared backpropagated signal at layer $\ell$. Under Kaiming
init, $q_\ell, \nu_\ell = \Theta(1)$ uniformly in $\ell$, so
\[
    \E\!\left[ \mathrm{tr}(F_\phip) \right]
    \;=\; \Theta\!\left( \sum_\ell M_\ell \right)
    \;=\; \Theta(\Cg).
\]

\paragraph{Step 3: From Fisher to Gauss-Newton.}
For softmax cross-entropy, the FIM and Gauss-Newton matrix differ by
a factor depending on the prediction distribution:
\[
    F_\phip \;=\; \E_X\!\left[ J_\phip^\top \cdot \mathrm{diag}(p)(I - 1 p^\top) \cdot J_\phip \right].
\]
At random initialization $p \to (1/N_{cls}) \mathbf{1}$, so the
inner matrix has eigenvalues bounded by $1/N_{cls}$. Therefore
\[
    \mathrm{tr}(F_\phip)
    \;\geq\; \frac{1}{N_{cls}^2} \cdot \mathrm{tr}(J_\phip^\top J_\phip)
\]
and combining with Step~2:
\[
    \E\!\left[ \mathrm{tr}(J_\phip^\top J_\phip) \right]
    \;=\; \Omega(\Cg).
\]

\paragraph{Step 4: Top eigenvalue from trace via participation ratio.}
The top eigenvalue is bounded from below by the trace divided by the
effective rank:
\[
    \lambda_{\max}(J_\phip^\top J_\phip)
    \;\geq\; \frac{\mathrm{tr}(J_\phip^\top J_\phip)}{\rho(J_\phip^\top J_\phip)},
\]
where $\rho$ is the participation ratio (effective rank under the
$L^2$ measure). Karakida et al.~\cite{karakida2019universal},
Theorem~4.2, show that under mean-field assumptions, the
participation ratio of the FIM is $\Theta(1)$ uniformly in network
width. Combining with Step~3:
\[
    \lambda_{\max}(J_\phip^\top J_\phip)
    \;=\; \Omega(\Cg).
\]

\paragraph{Step 5: Concentration in the wide-network limit.}
Karakida et al.~\cite{karakida2019universal}, Theorem~5.1, establish
that in the mean-field limit, the top eigenvalue of the FIM
concentrates around its expectation with high probability. The
deterministic bound therefore holds with probability $1 - o(1)$ as
$\Cg \to \infty$.

This completes the proof of Lemma~\ref{lem:phi_scaling}. \hfill$\square$

\begin{remark}[Empirical validation]
Karakida's asymptotic law $\lambda_{\max} = \Theta(M)$ corresponds
to a log-log scaling exponent of $0.5$ against $\Cg$ under
$\Cg = \Theta(M^2)$. Experiment~1.1 measures an empirical exponent of
$\approx 0.70$ over $\Cg$ spanning three decades
($D \in \{16, \ldots, 8192\}$), consistent with the asymptotic
$\sqrt{\Cg}$ prediction plus a mild finite-width upward correction
of the same character observed in Karakida's own numerical validation
(their Figures~2--3, exponent $\in [0.5, 0.8]$ across finite widths).
The qualitative claim of Lemma~\ref{lem:phi_scaling}---$\lambda_\phi$
grows without bound as the spatial module widens---is preserved at
all widths we test, and the quantitative slope matches the corrected
asymptotic prediction.
\end{remark}

\subsubsection{Proof of Lemma~\ref{lem:interlacing}}

The claim that $\lambda_{\max}(M) \geq \max\{\lambda_{\max}(A),
\lambda_{\max}(D)\}$ for a symmetric block matrix
$M = \begin{pmatrix} A & B \\ B^\top & D \end{pmatrix}$ is a direct
application of Cauchy's interlacing theorem to principal submatrices.
See Bhatia~\cite{bhatia1997matrix}, Corollary~III.1.5.

For the eigenvector $v$ achieving $\lambda_{\max}(A) = v^\top A v$,
the extension $\tilde{v} = \begin{pmatrix} v \\ 0 \end{pmatrix}$ has
$\tilde{v}^\top M \tilde{v} = v^\top A v$, so
$\lambda_{\max}(M) \geq \tilde{v}^\top M \tilde{v} / \norm{\tilde{v}}^2
= \lambda_{\max}(A)$. The same argument for $D$ completes the
proof. \hfill$\square$

\subsubsection{Proof of Lemma~\ref{lem:schur}}

We establish two claims. Part (i) is a rank bound derived from the
softmax simplex constraint, aggregated over the pixel grid and
dataset. Part (ii) is an upper bound on the smallest nonzero
eigenvalue via a trace-over-rank argument, using
Assumption~\ref{ass:inputlip} to bound the trace independently of
$\Cg$.

\paragraph{Convention on $\lambda_{\min}$.} The spectral block
$J_\thetap^\top J_\thetap$ is rank-deficient whenever the residual
dimension $N_{\mathrm{data}} HW (N_{\mathrm{cls}} - 1) < \Cf$
(which is the over-parameterised regime relevant for large spatial
classifiers). We interpret the bound on
$\lambda_{\min}(J_\thetap^\top J_\thetap)$ throughout this proof as a
bound on the smallest \emph{nonzero} eigenvalue $\lambda_{\min}^{+}$;
the strict smallest eigenvalue is zero and any positive upper bound
holds trivially. The condition number $\kappa$ appearing in
Theorem~\ref{thm:hessian} is understood as the effective condition
number $\kappa_{\mathrm{eff}}(G) = \lambda_{\max}(G) /
\lambda_{\min}^{+}(G)$ per
Remark~\ref{rem:kappa_eff} of Section~\ref{sec:setup}.

\paragraph{Step 1: Per-pixel softmax null-space constraint.}
For each pixel $i \in \{1, \dots, HW\}$ and each training example
$n \in \{1, \dots, N_{\mathrm{data}}\}$, the softmax output
$p_{n,i}(\thetap, \phip) \in \Delta^{N_{\mathrm{cls}} - 1}$ satisfies
$\mathbf{1}^\top p_{n,i} \equiv 1$ identically in $\thetap$.
Differentiating this identity in $\thetap$,
\[
    \mathbf{1}^\top \frac{\partial p_{n,i}}{\partial \thetap} \;=\; \mathbf{0}.
\]
The one-hot target $y_{n,i}$ also satisfies $\mathbf{1}^\top y_{n,i} =
1$, so the residual $r_{n,i} := p_{n,i} - y_{n,i}$ inherits the same
null-space constraint: $\mathbf{1}^\top r_{n,i} \equiv 0$ and
$\mathbf{1}^\top \partial r_{n,i} / \partial \thetap = \mathbf{0}$.
Consequently the per-pixel residual Jacobian block $\partial r_{n,i}
/ \partial \thetap \in \mathbb{R}^{N_{\mathrm{cls}} \times \Cf}$ has
row-rank at most $N_{\mathrm{cls}} - 1$: one linear dependency
between its $N_{\mathrm{cls}}$ rows is guaranteed.

\paragraph{Step 2: Aggregation over pixels and dataset.}
Stack the $N_{\mathrm{data}} HW$ per-pixel Jacobian blocks into
\[
    J_\thetap \;=\; \frac{\partial r}{\partial \thetap}
    \;\in\; \mathbb{R}^{N_{\mathrm{data}} HW N_{\mathrm{cls}} \times \Cf}.
\]
Each of the $N_{\mathrm{data}} HW$ blocks contributes at most
$N_{\mathrm{cls}} - 1$ independent rows (Step 1), and no cross-block
dependencies are imposed by the softmax simplex constraint. Hence
\[
    \mathrm{rank}(J_\thetap)
    \;\leq\; N_{\mathrm{data}} \cdot HW \cdot (N_{\mathrm{cls}} - 1).
\]
Combining with the trivial column-rank bound
$\mathrm{rank}(J_\thetap) \leq \Cf$ and the identity
$\mathrm{rank}(J_\thetap^\top J_\thetap) = \mathrm{rank}(J_\thetap)$,
we obtain claim (i) of the lemma:
\begin{equation}
    \mathrm{rank}(J_\thetap^\top J_\thetap)
    \;\leq\; \min\!\left(\Cf,\; N_{\mathrm{data}} HW (N_{\mathrm{cls}} - 1)\right).
    \label{eq:rank_ceiling}
\end{equation}
Denote the effective rank ceiling by $R := N_{\mathrm{data}} HW
(N_{\mathrm{cls}} - 1)$, which is $\Cg$-independent.

\paragraph{Step 3: Trace bound via encoder-bottleneck chain rule.}
Under Assumption~\ref{ass:linear}, $\partial Z_i / \partial \thetap =
I_K \otimes x_i^\top$
(Magnus and Neudecker~\cite{magnus1988matrix}, Ch.~9). Under
Assumption~\ref{ass:inputlip}, $\|\partial \hat{y}_i / \partial
Z_i\|_{\mathrm{op}} \leq \widetilde{L}$ uniformly. For the softmax
residual Jacobian at pixel $(n, i)$,
\[
    \frac{\partial r_{n,i}}{\partial \thetap}
    \;=\; \bigl[\mathrm{diag}(p_{n,i}) - p_{n,i} p_{n,i}^\top\bigr]
    \cdot \frac{\partial \hat{y}_{n,i}}{\partial Z_{n,i}}
    \cdot \frac{\partial Z_{n,i}}{\partial \thetap},
\]
with $\|\mathrm{diag}(p) - p p^\top\|_{\mathrm{op}} \leq 1$ (the softmax
Hessian has operator norm at most $1$; see Botev, Ritter, and
Barber~\cite{botev2017practical}). Taking the Frobenius norm,
\[
    \left\| \frac{\partial r_{n,i}}{\partial \thetap} \right\|_F^2
    \;\leq\; \widetilde{L}^2 \cdot \|I_K \otimes x_{n,i}^\top\|_F^2
    \;=\; \widetilde{L}^2 \cdot K \cdot \|x_{n,i}\|^2.
\]
Aggregating,
\begin{equation}
    \mathrm{tr}(J_\thetap^\top J_\thetap)
    \;=\; \sum_{n, i} \left\| \frac{\partial r_{n,i}}{\partial \thetap}
                       \right\|_F^2
    \;\leq\; \widetilde{L}^2 K \cdot \sum_{n, i} \|x_{n,i}\|^2
    \;\leq\; \widetilde{L}^2 K \cdot N_{\mathrm{data}} HW \cdot
             \mathrm{tr}(\Sigma_X),
    \label{eq:trace_bound}
\end{equation}
where $\Sigma_X$ is bounded by Assumption~\ref{ass:subg}. All factors
on the right-hand side are $\Cg$-independent.

\paragraph{Step 4: Trace/rank upper bound on $\lambda_{\min}^{+}$.}
For any positive semidefinite matrix $A$ of rank $r$ with eigenvalues
$0 = \lambda_1 = \ldots = \lambda_{n-r} < \lambda_{n-r+1} \leq \ldots
\leq \lambda_n$, the arithmetic-mean inequality on the nonzero
eigenvalues gives $\lambda_{n-r+1} \leq \mathrm{tr}(A) / r$.
Applied to $J_\thetap^\top J_\thetap$ with $r = \mathrm{rank}$ from
\eqref{eq:rank_ceiling} and trace bound \eqref{eq:trace_bound},
\[
    \lambda_{\min}^{+}(J_\thetap^\top J_\thetap)
    \;\leq\; \frac{\mathrm{tr}(J_\thetap^\top J_\thetap)}{\mathrm{rank}(J_\thetap^\top J_\thetap)}
    \;\leq\; \frac{\widetilde{L}^2 K \cdot N_{\mathrm{data}} HW \cdot
             \mathrm{tr}(\Sigma_X)}
             {N_{\mathrm{data}} HW (N_{\mathrm{cls}} - 1)}
    \;=\; \frac{\widetilde{L}^2 K \cdot \mathrm{tr}(\Sigma_X)}
                {N_{\mathrm{cls}} - 1}.
\]
Setting $C_2 := \widetilde{L}^2 K \cdot \mathrm{tr}(\Sigma_X)$ yields
claim (ii) of the lemma. The constant $C_2$ depends on the input
bottleneck dimension $K$, the classifier's depth-dependent input
Lipschitz constant $\widetilde{L}$, and the trace of the data
covariance $\mathrm{tr}(\Sigma_X)$---but not on $\Cg$, the width $M$,
or the spatial parameter count. This completes the proof.
\hfill$\square$

\begin{remark}[Corollary for per-pixel classifiers]
\label{rem:kron_corollary}
When $g_\phip$ is per-pixel (no spatial mixing), the per-pixel
classifier Jacobians $A_i \equiv A$ coincide, and the aggregated block
factors as $J_\thetap^\top J_\thetap = A^\top A \otimes \Sigma_X$ (per
example). Bhatia's Kronecker spectral
identity~\cite{bhatia1997matrix} then tightens the smallest positive
eigenvalue to
\[
    \lambda_{\min}^{+}(J_\thetap^\top J_\thetap)
    \;=\; \mu_r(A^\top A) \cdot \sigma_{\min}^{+}(\Sigma_X),
\]
where $\mu_r$ is the smallest nonzero eigenvalue of $A^\top A$. Both
factors are $\Cg$-independent. This tighter Corollary is what our
synthetic MLP experiments verify (Experiment~1.1); the general
Lemma~\ref{lem:schur} bound is what applies to real CNN/ViT spatial
classifiers.
\end{remark}

\begin{remark}[Effective rank vs.\ rank ceiling]
\label{rem:effective_rank}
Lemma~\ref{lem:schur}(i) is a bound on the rank ceiling, not on the
effective rank. Neural collapse~\cite{papyan2020prevalence} predicts
that in the terminal training phase only $\approx N_{\mathrm{cls}} - 1$
of the $N_{\mathrm{data}} HW (N_{\mathrm{cls}} - 1)$ residual
directions carry non-negligible curvature, so the effective rank
is much smaller than the ceiling. Our bound is therefore loose as a
descriptor of curvature-carrying directions but tight as the rank
ceiling; either interpretation suffices for
Theorem~\ref{thm:hessian}, since both imply
$\lambda_{\min}^{+}(J_\thetap^\top J_\thetap)$ is $\Cg$-independent.
\end{remark}

\begin{remark}[Softmax vs.\ sigmoid outputs]
The rank bound $HW(N_{\mathrm{cls}} - 1)$ depends on the simplex
constraint $\mathbf{1}^\top p = 1$, which holds for softmax and is
preserved under label smoothing. For sigmoid / independent-Bernoulli
multi-label losses the bound loosens to $HW N_{\mathrm{cls}}$
(no simplex constraint). Our spectroscopic pipelines use softmax
throughout, so the tighter bound applies.
\end{remark}

\subsubsection{Combining the three lemmas: full proof of Theorem~\ref{thm:hessian}}
\label{sec:thm1_combine}

Assemble Lemmas~\ref{lem:phi_scaling}, \ref{lem:interlacing}, and
\ref{lem:schur} to complete the proof of Theorem~\ref{thm:hessian}.
Recall the joint Gauss-Newton matrix
\[
    G \;=\; \begin{pmatrix}
        J_\thetap^\top J_\thetap & J_\thetap^\top J_\phip \\
        J_\phip^\top J_\thetap & J_\phip^\top J_\phip
    \end{pmatrix}
\]
and the effective condition number
$\kappa_{\mathrm{eff}}(G) = \lambda_{\max}(G) /
\lambda_{\min}^{+}(G)$ (Remark~\ref{rem:kappa_eff}).

\paragraph{Upper bound on $\lambda_{\max}(G)$ direction, from the spatial block.}
Lemma~\ref{lem:phi_scaling} establishes
\[
    \E\!\left[ \lambda_{\max}(J_\phip^\top J_\phip) \right]
    \;=\; \Omega(M)
    \;=\; \Omega\!\left(\sqrt{\Cg}\right)
\]
under Assumptions~\ref{ass:reg}--\ref{ass:subg} for fixed-depth
architectures with $\Cg = \Theta(M^2)$, with concentration in the
wide-network limit
(Karakida et al.~\cite{karakida2019universal}, Theorem~5.1).
Lemma~\ref{lem:interlacing} (Cauchy interlacing applied to the
diagonal block $J_\phip^\top J_\phip$ of $G$) transfers this to the
full matrix:
\[
    \lambda_{\max}(G)
    \;\geq\; \lambda_{\max}(J_\phip^\top J_\phip)
    \;=\; \Omega\!\left(\sqrt{\Cg}\right).
\]
Denote the implicit constant by $c_1'$: $\lambda_{\max}(G) \geq c_1'
\cdot \sqrt{\Cg}$ for all $\Cg$ sufficiently large.

\paragraph{Upper bound on $\lambda_{\min}^{+}(G)$ direction, from the spectral block.}
Lemma~\ref{lem:schur} (Step 6 of its proof) gives, under
Assumption~\ref{ass:linear},
\[
    \lambda_{\min}^{+}(J_\thetap^\top J_\thetap)
    \;\leq\; L^2 \cdot \sigma_{S}(\Sigma_X),
\]
where $L$ is the Lipschitz constant of $g_\phip$'s gradient
(Assumption~\ref{ass:reg}) and $\sigma_S(\Sigma_X) =
\lambda_{\min}(\Sigma_X)$ is the smallest data covariance eigenvalue.
By Cauchy interlacing (min form, Bhatia~\cite{bhatia1997matrix}
Corollary~III.1.5, dually):
\[
    \lambda_{\min}^{+}(G)
    \;\leq\; \lambda_{\min}^{+}(J_\thetap^\top J_\thetap)
    \;\leq\; L^2 \cdot \sigma_S(\Sigma_X).
\]
Denote this upper bound by $C_2 := L^2 \cdot \sigma_S(\Sigma_X)$.
Crucially, $C_2$ depends only on the model class of $g_\phip$
(via $L$) and on the data distribution (via $\sigma_S$)---it does
NOT depend on $\Cg$ or on the number of spatial parameters in any way.

\paragraph{Combining the two bounds.}
Dividing,
\begin{equation}
    \kappa_{\mathrm{eff}}(G)
    \;=\; \frac{\lambda_{\max}(G)}{\lambda_{\min}^{+}(G)}
    \;\geq\; \frac{c_1' \cdot \sqrt{\Cg}}{C_2}
    \;=\; \frac{c_1'}{C_2} \cdot \sqrt{\Cg}.
    \label{eq:kappa_bound_raw}
\end{equation}
This is the essential result: the effective condition number of the
joint Gauss-Newton matrix grows at least as $\sqrt{\Cg}$
(equivalently, linearly in the spatial module's characteristic
width $M$), with a rate constant $c_1' / C_2$ determined entirely by
the classifier's regularity and the data statistics.

\paragraph{Recasting in the $\sqrt{\Cg / \Cf}$ form of Theorem~\ref{thm:hessian}.}
Assumption~\ref{ass:linear} sets $\Cf = S \cdot K$, both fixed
architectural constants of the spectroscopic pipeline. Multiplying and
dividing \eqref{eq:kappa_bound_raw} by $\sqrt{\Cf}$:
\[
    \kappa_{\mathrm{eff}}(G)
    \;\geq\; \frac{c_1'}{C_2} \cdot \sqrt{\Cg}
    \;=\; \left( \frac{c_1' \cdot \sqrt{\Cf}}{C_2} \right) \cdot
                            \sqrt{\frac{\Cg}{\Cf}}
    \;=\; c_1 \cdot \sqrt{\frac{\Cg}{\Cf}},
\]
where $c_1 := c_1' \cdot \sqrt{\Cf} / C_2 = c_1' \cdot \sqrt{SK} / (L^2
\sigma_S(\Sigma_X))$ is the theorem's leading constant, depending only
on the data covariance ($\sigma_S$), the bottleneck dimension $K$,
the input spectral dimension $S$, and the model class of $g_\phip$
(through $L$)---exactly the dependencies claimed in the theorem
statement. The subtracted constant $c_2$ absorbs finite-$\Cg$
corrections from the mean-field concentration bound in
Lemma~\ref{lem:phi_scaling}. This completes the proof.
\hfill$\square$

\begin{remark}[Width vs.\ parameter-count scaling]
The $\sqrt{\Cg}$ scaling reflects the fact that Karakida's underlying
result~\cite{karakida2019universal} bounds $\lambda_{\max}$ in terms
of the network's characteristic \emph{width} $M$, not its total
parameter count. For a fixed-depth architecture with $\Cg =
\Theta(M^2)$, the equivalent statement is $\lambda_{\max} \sim
\sqrt{\Cg}$. The paper's Experiment~1.1, which sweeps width $D$ from
$16$ to $8192$ at fixed depth, measures a log--log slope of $\approx
0.70$ against $\Cg$---consistent with the asymptotic $0.5$
($\sqrt{\Cg}$) prediction plus a mild finite-width upward correction
of the same character reported in Karakida's own numerical
validation.
\end{remark}

\begin{remark}[Why the $\sqrt{\Cg / \Cf}$ framing]
Although $\Cf = SK$ is a constant of the setup and the bound
$\kappa_{\mathrm{eff}}(G) \geq \Omega(\sqrt{\Cg})$ is equivalent to
$\Omega(\sqrt{\Cg / \Cf})$ up to constants, we prefer the ratio form
because it makes explicit that the ill-conditioning is fundamentally
a \emph{capacity asymmetry} phenomenon: doubling both $\Cg$ and
$\Cf$ (e.g.\ by widening both modules together) leaves the pathology
unchanged. It is the ratio, not the absolute size, that drives the
theorem.
\end{remark}

\subsection{Full proof of Theorem~\ref{thm:twoscale}}
\label{sec:proof_thm2}

\TODO{Phase 3, weeks 8-10. Combine Borkar singular perturbation with
Pezeshki residual decay. Proof skeleton: (i) ODE limits of SGD,
(ii) timescale gap from Lemma 5.1, (iii) Tikhonov fast-manifold
convergence, (iv) NTK residual decay along fast manifold,
(v) slow-variable bound by Gronwall.}

\subsection{Empirical robustness}

We here report robustness experiments referenced in the main text.

\subsubsection{Optimizer ablation (SGD vs Adam)}
\TODO{Phase 5 follow-up. Show EGR collapse pattern in both regimes;
note Adam masks the natural gradient ratio.}

\subsubsection{Loss-variant ablation}
\TODO{Phase 5 follow-up. Cross-entropy, focal loss, label smoothing.}

\subsubsection{Regularizer ablation}
\TODO{Phase 5 follow-up. Weight decay, dropout, batch norm.
Pezeshki et al.~\cite{pezeshki2021gradient}, App.~B, already show
gradient starvation is robust to these; we reproduce briefly in our
setting.}

\subsubsection{Smallest capacity ratio where the pathology disappears}
\TODO{Phase 5 follow-up. Sweep $\Cg/\Cf$ from 1 to 1000, identify
the threshold below which joint training reaches frozen
performance.}
